\chapter{Source Code Walkthrough}
\label{chap:source_tour}

The defining hallmark of Andrew Tanenbaum's \textit{Operating Systems: Design and Implementation} was the complete printed source code of MINIX. Students did not simply study abstract concepts; they examined the exact C structs, function signatures, and assembly stubs that brought the operating system to life.

In this chapter, we provide an annotated guided tour of the **Hydra operating system source tree**. We trace the life of the system from UEFI initialization to Ring~0 takeover, examine the key modules in \texttt{src/}, and dissect the Ring~3 isolated userland drivers in \texttt{userland/}.

% ==============================================================================
\section{Directory Layout and Crate Organization}
\label{sec:code_layout}
% ==============================================================================

Hydra is organized into clean, modular Rust source files structured across two primary tiers: the Ring~0 microkernel core (\texttt{src/}) and the Ring~3 isolated driver servers (\texttt{userland/drivers/}):

\begin{verbatim}
d:\os
|-- src/
|   |-- arch/                 # x86_64 CPU, GDT, IDT, PIC, APIC, MSR syscall stubs
|   |-- capability/           # Capability-based security tokens and rights tables
|   |-- compat/linux/         # Linux ABI syscall dispatcher, POSIX VFS, devices
|   |-- fs/zfs/               # OpenZFS storage engine, uberblocks, space maps, dnodes
|   |-- ipc/                  # Synchronous rendezvous endpoints, shared memory
|   |-- memory/               # Physical frame allocator, 4-level PML4 paging, heap
|   |-- net/                  # Network stack: ARP, IPv4, ICMP, UDP, TCP state machine
|   |-- task/                 # Task Control Block (TCB), preemptive scheduler, ELF
|   |-- video/                # UEFI GOP framebuffer and DRM/KMS card0 emulation
|   +-- main.rs               # UEFI entry point and milestone boot orchestrator
|-- userland/
|   |-- drivers/
|   |   |-- ata-driver/       # Ring 3 IDE/ATA Bus-Master DMA storage server (0x21)
|   |   |-- zfs-server/       # Ring 3 OpenZFS storage daemon (0x20)
|   |   |-- e1000-driver/     # Ring 3 Intel 82540EM Gigabit NIC server (0x22)
|   |   +-- display-server/   # Ring 3 DRM/KMS graphics compositor server (0x23)
|   +-- hydra-logind/         # D-Bus session and seat management daemon
+-- book/                     # The Hydra: Design and Implementation textbook
\end{verbatim}

% ==============================================================================
\section{The Microkernel Entry Point (\texttt{src/main.rs})}
\label{sec:tour_main}
% ==============================================================================

Hydra begins execution in \texttt{src/main.rs} via the standard UEFI entry point:
\begin{lstlisting}[language=Rust, caption={The Hydra Entry Point (\texttt{src/main.rs})}]
#[no_mangle]
pub extern "efiapi" fn efi_main(image_handle: Handle, system_table: SystemTable<Boot>) -> Status {
    // Step 1: Initialize early serial logging (COM1 at 0x3F8)
    arch::raw_serial::init(0x3F8);
    kprintln!("[*] Step 1: Hydra Microkernel Initializing...");

    // Step 2: Initialize UEFI Graphics Output Protocol (GOP)
    let fb_info = video::init_gop(&system_table)?;

    // Step 3: Capture memory map and execute ExitBootServices
    arch::boot_services::take_hardware_ownership();

    // Step 4: Configure x86_64 GDT, TSS, IDT, and Interrupt Gates
    arch::gdt::init();
    arch::idt::init();

    // Step 5: Initialize 4-Level Paging and Frame Allocator
    memory::init();

    // Step 6: Enable MSR SYSCALL/SYSRET Dispatcher
    arch::syscall::init_msrs();

    // Step 7: Mount OpenZFS Pool ("hydrapool")
    fs::zfs::init_pool("hydrapool")?;

    // Step 8: Spawn Ring 3 System Servers
    supervisor::spawn_server("ata_driver", "/bin/ata_driver.elf", EndpointId(0x21));
    supervisor::spawn_server("zfs_server", "/bin/zfs_server.elf", EndpointId(0x20));
    supervisor::spawn_server("net_driver", "/bin/net_driver.elf", EndpointId(0x22));
    supervisor::spawn_server("display_server", "/bin/display_server.elf", EndpointId(0x23));

    // Step 9: Launch Linux Chroot Environment and KDE Session
    task::scheduler::start_multitasking();
}
\end{lstlisting}

% ==============================================================================
\section{System Call Interception (\texttt{src/compat/linux/mod.rs})}
\label{sec:tour_syscalls}
% ==============================================================================

When a user process executes \texttt{syscall}, the assembly entry routine captures the register frame and calls \texttt{linux\_syscall\_dispatcher}:
\begin{lstlisting}[language=Rust, caption={Linux System Call Multiplexing in Hydra}]
pub fn linux_syscall_dispatcher(frame: &mut SyscallFrame) -> i64 {
    let nr = frame.rax;
    match nr {
        0 => sys_read(frame.rdi as i32, frame.rsi, frame.rdx as usize),
        1 => sys_write(frame.rdi as i32, frame.rsi, frame.rdx as usize),
        2 => sys_open(frame.rdi, frame.rsi as i32, frame.rdx as u32),
        3 => sys_close(frame.rdi as i32),
        9 => sys_mmap(frame.rdi, frame.rsi, frame.rdx, frame.r10, frame.r8, frame.r9),
        56 => sys_clone(frame),
        59 => sys_execve(frame),
        232 => sys_epoll_wait(frame.rdi as i32, frame.rsi, frame.rdx as i32, frame.r10 as i32),
        _ => {
            kprintln!("[linux] unhandled syscall {} -> ENOSYS", nr);
            -ENOSYS
        }
    }
}
\end{lstlisting}

% ==============================================================================
\section{The Ring~3 Driver Loop: An Example}
\label{sec:tour_drivers}
% ==============================================================================

How does a driver server execute in Ring~3? The ATA disk driver (\texttt{userland/drivers/ata-driver/src/main.rs}) demonstrates the standard server event loop:
\begin{lstlisting}[language=Rust, caption={The Ring 3 Driver Event Loop}]
fn main() {
    let endpoint = Endpoint::bind(0x21).expect("Failed to bind ATA endpoint");
    let mut msg_buffer = [0u8; 1024];

    loop {
        // Block waiting for client request (synchronous rendezvous)
        let (client_id, request_type) = endpoint.receive(&mut msg_buffer);

        match request_type {
            CMD_READ_SECTORS => {
                let lba = parse_lba(&msg_buffer);
                let count = parse_count(&msg_buffer);
                dma_read_sectors(lba, count);
                endpoint.reply(client_id, STATUS_OK, &[]);
            }
            CMD_WRITE_SECTORS => {
                let lba = parse_lba(&msg_buffer);
                dma_write_sectors(lba, &msg_buffer);
                endpoint.reply(client_id, STATUS_OK, &[]);
            }
            _ => {
                endpoint.reply(client_id, STATUS_INVALID_CMD, &[]);
            }
        }
    }
}
\end{lstlisting}

% ==============================================================================
\section{Summary}
\label{sec:tour_summary}
% ==============================================================================

In this chapter, we walked through the living codebase of Hydra:
\begin{itemize}
    \item \texttt{src/main.rs} cleanly coordinates the milestone boot sequence from raw UEFI to preemptive multitasking.
    \item System calls enter through hardware MSRs and demultiplex into strongly typed Rust handler functions.
    \item Ring~3 isolated driver processes operate on clean request/reply event loops over synchronous IPC endpoints.
    \item The complete code reflects the vision of Tanenbaum: modular, clean, understandable, and engineered for high reliability.
\end{itemize}
